\documentclass[11pt]{article}
\usepackage[T1]{fontenc}
\usepackage[utf8]{inputenc}
\usepackage[margin=1in]{geometry}
\usepackage{parskip}
\usepackage{hyperref}
\pagestyle{empty}
\setlength{\parskip}{0.6em}

\begin{document}

{\noindent Dietmar Henrich \,|\, Chair of Materials Science and Nanotechnology,
Technische Universit\"{a}t Dresden, 01062 Dresden, Germany \,|\,
\href{mailto:profhenrich@googlemail.com}{profhenrich@googlemail.com}}

\vspace{0.8em}
\noindent \today

\vspace{0.8em}
\noindent The Editors, \emph{American Journal of Physics}

\noindent Dear Editors,

We are pleased to submit our manuscript, \textbf{``Memory as Dissipation: A
Curvature-Adaptive Dynamical Solver for the Lagrange Points of the Restricted
Three-Body Problem, with an N-body Stability Cross-Check,''} for consideration as
a Paper in the \emph{American Journal of Physics}.

The manuscript revisits a problem familiar to every student of mechanics---the
five Lagrange points of the restricted three-body problem---but uses it as a
transparent laboratory for a conceptual question that we believe is well suited
to your readership: \emph{when a physical dynamical system is said to ``compute''
the solution of a problem, which term in its equations of motion actually does
the work?} We make the answer concrete and, we hope, instructive. Building a
solver in the style of physics-based (memcomputing) approaches, we first show by
an elementary kinetic-energy identity that the intuitive choice---a growing
memory variable multiplying the gradient force---is \emph{dynamically inert}: a
multiplier on a conservative force can only exchange energy with the potential
and can never dissipate it. We then show that folding the same memory into the
\emph{dissipation} makes it genuinely load-bearing, and we verify every claim
numerically. The argument requires no machinery beyond Newtonian mechanics, the
work--energy theorem, and a simple numerical integrator, yet it yields a lesson
that we think students and instructors will find memorable: \emph{check the
energy budget before believing that a variable computes}.

We believe the paper fits the mission of \emph{AJP} for several reasons. (i)~It
treats a classic, curriculum-central problem with a fresh conceptual angle rather
than new apparatus. (ii)~Its derivations are self-contained and accessible at the
advanced-undergraduate / beginning-graduate level. (iii)~It is honest about
scope: we state plainly that classical root-finders are faster on this particular
problem, that the method's value here is pedagogical and diagnostic, and we
characterise---rather than hide---a failure mode at extreme mass ratios.
(iv)~The work is fully reproducible: an open-source repository with interactive
browser applications lets readers and students vary the parameters and watch the
memory, the dissipation it generates, and the kinetic energy decay in real time,
and a self-contained N-body module reproduces the Trojan-stability cross-check
(the measured Jupiter libration period agrees with theory to $0.1\%$).

The manuscript is original, has not been published previously, and is not under
consideration by any other journal. All authors have approved the submission and
declare no conflicts of interest. The corresponding author is Dietmar Henrich.

Thank you for considering our work. We would be glad to respond to any questions.

\vspace{0.8em}
\noindent Sincerely,\\[0.4em]
Dietmar Henrich \quad$\cdot$\quad Massimiliano Di Ventra

\end{document}
